\documentclass[11pt,a4paper]{article}

% ---------------------------------------------------------
% PACKAGES
% ---------------------------------------------------------
\usepackage[utf8]{inputenc}
\usepackage[T1]{fontenc}
\usepackage{lmodern}
\usepackage{amsmath, amssymb, amsfonts, bm}
\usepackage{geometry}
\usepackage{graphicx}
\usepackage{hyperref}
\usepackage{authblk}
\usepackage{cite}
\usepackage{physics}
\usepackage{mathtools}
\usepackage{tensor}
\usepackage{enumitem}

\geometry{margin=2.4cm}

% ---------------------------------------------------------
% TITLE
% ---------------------------------------------------------
\title{\textbf{Companion VIII (v49.1)}\\[0.2cm]
\textbf{Likelihood Structure and Global Survey Interface}\\[0.2cm]
\large Foundations of Relaxation-Driven Cyclic Cosmology (RDCC)}

\author[1]{Michael Lehmann}
\affil[1]{\small Independent Researcher, Relaxation-Driven Cyclic Cosmology (RDCC)\\
\texttt{mi.lehmann@gmx.de}}

\date{September 2026}

\begin{document}
\maketitle

% ========================================================
% ABSTRACT
% ========================================================

\begin{abstract}
This Companion develops the corrected likelihood structure of RDCC in its
Phase-1-integrated v49.1 form. The global survey interface combines CMB, BAO,
RSD multipoles, supernovae, and gravitational-wave observables into a single
covariant likelihood. The covariance matrix encodes cross-correlations between
visible and mirror-sector contributions and incorporates the dimension-8
inter-sector operator $T^{(+)}_{\mu\nu}T^{(-)\mu\nu}$. The coefficients
$C_1$–$C_4$ are fixed theoretical factors of order unity rather than fit
parameters. This Companion provides the statistical backbone of RDCC and
connects the theoretical framework to observational data.
\end{abstract}

\tableofcontents
\newpage
% ========================================================
% 1 INTRODUCTION
% ========================================================

\section{Introduction}

RDCC is a global framework. Its observational interface must therefore be
constructed in a way that respects the bipartite Hilbert space structure and
the global CPT symmetry. The likelihood must combine all relevant probes:
\begin{itemize}[leftmargin=0.8cm]
  \item CMB temperature and polarization,
  \item BAO distance measurements,
  \item RSD multipoles,
  \item type Ia supernovae,
  \item gravitational-wave observables.
\end{itemize}

In v49.1, the likelihood structure is corrected to:
\begin{itemize}[leftmargin=0.8cm]
  \item treat $C_1$–$C_4$ as fixed theoretical factors,
  \item incorporate the dimension-8 inter-sector operator,
  \item encode cross-correlations between visible and mirror-sector
        contributions,
  \item remain fully EFT-coherent and statistically consistent.
\end{itemize}

This Companion develops the global likelihood structure and its connection to
the RDCC parameter space.
% ========================================================
% 2 GLOBAL DATA VECTOR
% ========================================================

\section[Global Data Vector]{Global Data Vector}

The global data vector is defined as
\begin{equation}
\mathbf{d}
=
\begin{pmatrix}
\mathbf{d}_{\mathrm{CMB}} \\
\mathbf{d}_{\mathrm{BAO}} \\
\mathbf{d}_{\mathrm{RSD}} \\
\mathbf{d}_{\mathrm{SN}} \\
\mathbf{d}_{\mathrm{GW}}
\end{pmatrix},
\end{equation}
where each block contains the corresponding observables:
\begin{itemize}[leftmargin=0.8cm]
  \item $\mathbf{d}_{\mathrm{CMB}}$: temperature and polarization spectra,
  \item $\mathbf{d}_{\mathrm{BAO}}$: distance indicators,
  \item $\mathbf{d}_{\mathrm{RSD}}$: multipoles of the power spectrum,
  \item $\mathbf{d}_{\mathrm{SN}}$: luminosity distances,
  \item $\mathbf{d}_{\mathrm{GW}}$: stochastic gravitational-wave background
        and standard sirens.
\end{itemize}

The theoretical prediction is
\begin{equation}
\mathbf{t}(\boldsymbol{\theta})
\end{equation}
with parameter vector $\boldsymbol{\theta}$ containing:
\begin{itemize}[leftmargin=0.8cm]
  \item standard cosmological parameters,
  \item RDCC-specific parameters (e.g. $\alpha_{\mathrm{IR}}$),
  \item survey nuisance parameters.
\end{itemize}
% ========================================================
% 3 COVARIANCE MATRIX
% ========================================================

\section[Covariance Matrix]{Covariance Matrix}

The covariance matrix is defined as
\begin{equation}
\mathbf{C}
=
\begin{pmatrix}
\mathbf{C}_{\mathrm{CMB}} & \mathbf{X}_{\mathrm{CMB-BAO}} & \cdots \\
\mathbf{X}_{\mathrm{BAO-CMB}} & \mathbf{C}_{\mathrm{BAO}} & \cdots \\
\vdots & \vdots & \ddots
\end{pmatrix},
\end{equation}
where:
\begin{itemize}[leftmargin=0.8cm]
  \item $\mathbf{C}_{\mathrm{CMB}}$ is the CMB covariance,
  \item $\mathbf{C}_{\mathrm{BAO}}$ is the BAO covariance,
  \item $\mathbf{C}_{\mathrm{RSD}}$ is the RSD covariance,
  \item $\mathbf{C}_{\mathrm{SN}}$ is the supernova covariance,
  \item $\mathbf{C}_{\mathrm{GW}}$ is the gravitational-wave covariance,
  \item $\mathbf{X}_{\mathrm{CMB-BAO}}$, etc., encode cross-correlations.
\end{itemize}

The dimension-8 operator contributes to the covariance through geometric
corrections,
\begin{equation}
\mathbf{C}
=
\mathbf{C}^{(0)}
+
\frac{C_2}{M^4} \, \mathbf{C}^{(\mathrm{int})},
\end{equation}
where $C_2$ is a fixed theoretical factor of order unity.
% ========================================================
% 4 LIKELIHOOD FUNCTION
% ========================================================

\section[Likelihood Function]{Likelihood Function}

The likelihood is defined as
\begin{equation}
\mathcal{L}(\boldsymbol{\theta})
=
\exp\left(
-\frac{1}{2}
\Delta\mathbf{d}^{\mathrm{T}}
\mathbf{C}^{-1}
\Delta\mathbf{d}
\right),
\end{equation}
with residual vector
\begin{equation}
\Delta\mathbf{d}
=
\mathbf{d}
-
\mathbf{t}(\boldsymbol{\theta}).
\end{equation}

The RDCC-specific parameters enter through:
\begin{itemize}[leftmargin=0.8cm]
  \item the modified Boltzmann hierarchy (Companion VII),
  \item the relaxation geometry (Companion VI),
  \item the tensor sector (Companion II and XI),
  \item the thermodynamic sector (Companion IV),
  \item the ultraviolet sector (Companion X).
\end{itemize}
% ========================================================
% 5 ROLE OF C-COEFFICIENTS
% ========================================================

\section[Role of C coefficients]{Role of $C_1$–$C_4$ Coefficients}

The coefficients $C_1$–$C_4$ appear in:
\begin{itemize}[leftmargin=0.8cm]
  \item the covariant Poisson equation (Companion VII),
  \item the covariance matrix corrections,
  \item the survey interface (Companion XII).
\end{itemize}

They satisfy
\begin{align}
C_1 &\sim \mathcal{O}(1), \\
C_2 &\sim \mathcal{O}(1), \\
C_3 &\sim \mathcal{O}(1), \\
C_4 &\sim \mathcal{O}(1),
\end{align}
and are fixed theoretical factors, not fit parameters. Their values arise from
the global geometric structure and the dimension-8 operator.
% ========================================================
% 6 CONCLUSION
% ========================================================

\section{Conclusion}

This Companion develops the corrected likelihood structure of RDCC in its
v49.1 form. The global survey interface combines CMB, BAO, RSD multipoles,
supernovae, and gravitational-wave observables into a single covariant
likelihood. The covariance matrix encodes cross-correlations between visible
and mirror-sector contributions and incorporates the dimension-8 inter-sector
operator.

The coefficients $C_1$–$C_4$ are fixed theoretical factors of order unity,
ensuring that the statistical structure is fully determined by the global
geometry. These results integrate the Boltzmann hierarchy (Companion VII), the
relaxation sector (Companion III and VI), the ultraviolet sector (Companion X),
and the survey interface (Companion XII). They form the backbone of RDCC’s
connection to observational data and parameter estimation (Companion IX).

\end{document}
